\chapter{System Architecture and Design}

\section{Overview of the System}

The API Development Tool follows a two-tier client-server model with three functional layers: a React workflow frontend, a Node.js/Express orchestration backend, and an enterprise connector repository containing module structure, scripts, and Maven tooling. Communication is split between REST (CRUD and control operations) and Socket.IO (real-time logs and sync events). The backend links web interactions to filesystem artifacts and external generation scripts. Figure~\ref{fig:system_arch} shows this layered boundary.

\begin{figure}[H]
\centering
\begin{tikzpicture}[
    comp/.style={rectangle, rounded corners=4pt, line width=0.6pt,
                 text width=2.4cm, align=center, minimum height=1cm,
                 font=\footnotesize, inner sep=5pt,
                 drop shadow={shadow xshift=0.5pt, shadow yshift=-0.5pt, opacity=0.15}},
    layer/.style={rectangle, rounded corners=8pt, draw=#1!45, line width=1pt,
                  inner xsep=10pt, inner ysep=20pt, fill=#1!6},
    layerlbl/.style={font=\footnotesize\bfseries, text=black!55},
    plbl/.style={font=\scriptsize, rounded corners=3pt, inner sep=3pt,
                 text=black!80, fill=white, draw=black!15},
    node distance=0.35cm
]

% --- Presentation Layer components ---
\node[comp, fill=blue!15, draw=blue!50!black] (yu)
    {\textbf{YamlUploader}\\{\scriptsize Workflow Engine}};
\node[comp, fill=blue!15, draw=blue!50!black, right=of yu] (cv)
    {\textbf{ConfigViewer}\\{\scriptsize Editor}};
\node[comp, fill=blue!15, draw=blue!50!black, right=of cv] (sp_c)
    {\textbf{Swagger Parser}\\{\scriptsize Client-Side}};
\node[comp, fill=blue!15, draw=blue!50!black, right=of sp_c] (jp_c)
    {\textbf{JAR Parser}\\{\scriptsize Client-Side}};

% --- Application Layer components ---
\node[comp, fill=green!14, draw=green!45!black, below=3.5cm of yu] (rest)
    {\textbf{REST API}\\{\scriptsize Routes}};
\node[comp, fill=green!14, draw=green!45!black, right=of rest] (sio)
    {\textbf{Socket.IO}\\{\scriptsize Event Handler}};
\node[comp, fill=green!14, draw=green!45!black, right=of sio] (gen)
    {\textbf{Generation}\\{\scriptsize Orchestrator}};
\node[comp, fill=green!14, draw=green!45!black, right=of gen] (sync)
    {\textbf{Sync Engine}\\{\scriptsize Chokidar}};

% --- Storage Layer components ---
\node[comp, fill=orange!16, draw=orange!50!black, below=3.5cm of rest,
      text width=3cm] (db)
    {\textbf{sql.js Database}\\{\scriptsize In-Process SQLite}};
\node[comp, fill=orange!16, draw=orange!50!black, right=0.5cm of db,
      text width=4.2cm] (fstore)
    {\textbf{Filesystem Storage}\\{\scriptsize swagger/ ~ configs/ ~ client-kits/}};

% --- External System ---
\node[comp, fill=red!14, draw=red!45!black, below=3cm of db,
      text width=5cm, xshift=1.8cm] (repo)
    {\textbf{Digital Connector Repository}\\{\scriptsize Module tree ~ | ~ Scripts ~ | ~ Maven build}};

% --- Layer backgrounds (drawn BEHIND components) ---
\begin{pgfonlayer}{background}
\node[layer=blue, fit=(yu)(cv)(sp_c)(jp_c)] (fe_layer) {};
\node[layer=green, fit=(rest)(sio)(gen)(sync)] (be_layer) {};
\node[layer=orange, fit=(db)(fstore)] (st_layer) {};
\node[layer=red, fit=(repo)] (ext_layer) {};
\end{pgfonlayer}

% --- Layer titles INSIDE each box (clear of the inter-layer gaps) ---
\node[layerlbl, anchor=north] at ([yshift=-3pt]fe_layer.north)
    {Presentation Layer --- React 18 Frontend (Tomcat WAR)};
\node[layerlbl, anchor=north] at ([yshift=-3pt]be_layer.north)
    {Application Layer --- Node.js / Express Backend (PM2)};
\node[layerlbl, anchor=north] at ([yshift=-3pt]st_layer.north)
    {Storage Layer};
\node[layerlbl, anchor=north] at ([yshift=-3pt]ext_layer.north)
    {External System};

% --- Protocol-coloured arrows (anchored to component x-positions) ---
\draw[<->, line width=1.1pt, >=stealth, color=blue!55!black]
    (cv.south |- fe_layer.south) --
    node[plbl, left, xshift=-4pt] {REST / HTTPS}
    (cv.south |- be_layer.north);
\draw[<->, line width=1.1pt, >=stealth, color=teal!65!black]
    (sp_c.south |- fe_layer.south) --
    node[plbl, right, xshift=4pt] {Socket.IO}
    (sp_c.south |- be_layer.north);
\draw[<->, line width=1.1pt, >=stealth, color=orange!55!black]
    (rest.south |- be_layer.south) --
    node[plbl, left, xshift=-4pt] {SQL Queries}
    (rest.south |- st_layer.north);
\draw[<->, line width=1.1pt, >=stealth, color=green!50!black]
    (gen.south |- be_layer.south) --
    node[plbl, right, xshift=4pt] {File I/O + Watchers}
    (gen.south |- st_layer.north);
\draw[->, line width=1.1pt, >=stealth, color=red!50!black]
    (db.south |- st_layer.south) --
    node[plbl, right, xshift=4pt] {PowerShell spawn}
    (db.south |- ext_layer.north);

\end{tikzpicture}
\caption{Layered system architecture of the API Development Tool, showing internal subcomponents, communication protocols, and deployment boundaries.}
\label{fig:system_arch}
\end{figure}

\section{Frontend Architecture}

The frontend is a React 18~\cite{react_docs} single-page app with React Router under the \texttt{/retail} base path. The root route hosts \texttt{YamlUploader} (primary workflow), while \texttt{/configs} hosts \texttt{ConfigViewer} (configuration review/editing).

\texttt{YamlUploader} centralizes workflow state for module context, specification data, JAR metadata, mappings, generation status, and console output. It establishes Socket.IO subscriptions and a polling fallback for environments with unstable WebSocket connectivity. Backend calls are wrapped in a unified API service with environment-driven base URL resolution. \texttt{ConfigViewer} reuses parser/generator utilities to keep editing behavior consistent with initial generation.

\section{Utility Layer: Swagger Parser}

The Swagger parser converts OpenAPI/Swagger documents into mapping-ready metadata through three linked steps: \texttt{\$ref} resolution, schema flattening, and operation extraction.

\texttt{\$ref} targets are resolved by traversing JSON pointer segments, with \texttt{WeakMap}-based caching to avoid repeated work. Flattening then produces deduplicated dot-path fields (including arrays and composed schemas via \texttt{allOf}/\texttt{oneOf}/\texttt{anyOf}) with circular-reference protection and depth bounds. Operation extraction combines these results to produce parameter and schema details per endpoint, including precedence-safe merge of path-level and operation-level parameters and support for both Swagger 2.0 and OpenAPI 3.0 request-body styles.

\section{Utility Layer: JAR Parser and Bytecode Introspection}

The JAR parser extracts EJB metadata directly from compiled archives without source code or JVM runtime requirements. Browser and server variants share the same core logic.

JAR files are read via \texttt{jszip}; class files are processed by \texttt{java-class-tools} in batches. The parser decodes JVM descriptors and generic signatures, resolves hierarchy details, and identifies service interfaces using \texttt{javax.ejb.Remote}/\texttt{jakarta.ejb.Remote} annotations plus naming heuristics.

Its critical output is deep, navigable field paths for method inputs and outputs. Recursive traversal includes inherited fields, getter-derived fields, and generic collection element types, bounded by configured depth and path caps for stability. Compiler-generated parameter names are normalized to meaningful names using type-based heuristics.

\section{Utility Layer: Configuration Generator}

The configuration generator converts interactive mappings into YAML required by the generation pipeline. It handles path parameters, request mappings, response mappings, and multi-EJB call ordering.

Request-side logic normalizes target paths (prefix cleanup, array notation, casing), supports source inference for unmapped fields, and rewrites interface-derived prefixes to implementation-compatible paths when needed. Response-side logic filters non-business fields, derives usable nested paths, and supports chained mappings across multiple EJB calls.

The output YAML includes module metadata, generation flags, endpoint definitions, EJB call sequences, and field correspondences, serialized through \texttt{js-yaml}.

The end-to-end transformation is linear and implementation-bound: Swagger/OpenAPI parsing and JAR introspection produce field metadata, interactive mapping resolves source-target correspondences, and the resulting YAML feeds module generation.

\section{Backend Architecture}

The backend is built on Express.js~\cite{express_docs} and combines routing, storage, synchronization, real-time messaging, and generation orchestration in one service. Startup initializes environment variables, storage paths, directories, sql.js schema/migrations, reconciliation routines, and HTTP/HTTPS server plus Socket.IO.

REST endpoints are grouped by domain: modules, Swagger files, client kits, user configs, and generation, with a health endpoint for diagnostics. Metadata persistence uses sql.js~\cite{sqljs_docs} tables (\texttt{swagger\_files}, \texttt{client\_kit\_metadata}, \texttt{user\_configs}) with uniqueness constraints and synchronous export-on-write durability. Validation logic enforces file type and path safety (including path traversal checks) and normalizes module identifiers.

\section{Filesystem Synchronisation}

Consistency between disk state and metadata is maintained in two phases: startup reconciliation and runtime monitoring. Startup routines re-import missing files, remove stale records, and repair drift for modules, client kits, and configs.

At runtime, chokidar~\cite{chokidar_docs} watchers monitor module, client-kit, and config paths with debounced processing and sync coalescing. Confirmed changes trigger database updates and Socket.IO broadcasts (\texttt{modules:sync}, \texttt{swaggers:sync}, \texttt{clientkit:sync}, \texttt{configs:sync}). Generation events (\texttt{generate:start}, \texttt{generate:output}, \texttt{generate:done}) provide live execution visibility.

\section{Module Generation Pipeline}

Generation converts saved Swagger and user-config artifacts into deployable module code. The backend locates the generation script, invokes it through \texttt{spawn}, and supports safe-add (default), force, and optional skip-build modes.

Process output is streamed line-by-line via \texttt{generate:output}, and completion status is emitted through \texttt{generate:done} after success/failure marker checks. A ten-minute timeout protects against hung executions. On platforms without PowerShell, the endpoint returns HTTP 501 while leaving all non-generation features available.

\section{Deployment Architecture}

Deployment uses two services: frontend WAR on Tomcat~\cite{tomcat_docs} and backend Node.js service under PM2. Packaging is automated by a ten-step PowerShell flow covering dependency checks, frontend build, WAR creation, optional smoke tests, and final ZIP bundling.

All runtime paths derive from a configurable digital root, allowing the same codebase to run on Windows and Linux without hardcoded environment-specific paths.

\section{Key Design Decisions}

The implementation prioritizes self-contained deployment, operational transparency, and compatibility with existing enterprise tooling. Filesystem-centric artifact storage plus sql.js export-on-write persistence offers simple auditability and durability for team-scale use.

Workflow state centralization in \texttt{YamlUploader} favors end-to-end usability, while PowerShell-based generation preserves compatibility with established Windows-centric pipelines. Avoiding external database dependencies keeps deployment viable in air-gapped environments.
